% Separate post-CP17 research. Full root/B/C/D actual-file audits PASS.
\subsection{Canonical module jets from two fixed parameter thickenings}

\begin{lemma}[Proved: uniform saturation loss for canonical pullbacks]
\label{D:canonical-pullback-uniform-jet-loss}
Let $R$ be a fixed complete excellent normal local surface domain
over an algebraically closed coefficient field $k$ of characteristic
$p>0$. Fix the ideal $J$ of
Propositions~\ref{D:cover-uniform-canonical-defect} and
\ref{D:cover-uniform-Noether-different}. Choose once a generically
separable finite Noether normalization
\[
 S=k[[x,y]]\subset R,\qquad x,y\in J,
 \qquad \mathfrak q=(x,y)\subset S.
                                                        \tag{D.CJ.1}
\]
There is an integer $c\ge0$, depending only on these fixed choices,
with the following property. Let $A/R$ be a finite normal cover domain
of degree $n$, etale on the puncture, let
$\pi:Y\to\operatorname{Spec}A$ be a projective regular resolution
which is an isomorphism on the puncture, and let $M$ be a reflexive
$A$-module. Set
\[
 F=(\pi^*M)^{\vee\vee},\qquad I_b=(x^b,y^b)A,
 \qquad K_b=H^0(Y,I_b\mathcal O_Y\cdot F)\subset M.
\]
Then, for $b\ge2c$,
\[
                 K_b\subset\mathfrak q^{\,b-2c}M.
                                                        \tag{D.CJ.2}
\]
In particular, if $\ell\ge1$ and $b\ge2\ell-1+2c$, then
$K_b\subset I_\ell M$. The constants do not depend on $A$, $M$,
its rank, or the resolution.
\end{lemma}

\begin{proof}
The normalization in (D.CJ.1) can be chosen inside $J$.
Start with a separating finite normalization $k[[u,v]]\subset R$.
For a sufficiently large $p$-power $q$, put
$x=u^q(1+u)$ and $y=v^q(1+v)$. These lie in $J$, generate an
$\mathfrak m_R$-primary ideal, and give a finite normalization.
The monic equations $u^{q+1}+u^q-x=0$ and
$v^{q+1}+v^q-y=0$ have nonzero derivatives $u^q,v^q$;
thus $\operatorname{Frac}k[[u,v]]/\operatorname{Frac}k[[x,y]]$
is separable, as is its composite with the original separating
normalization. One may add this normalization to the fixed finite
list used in the cited Noether-different proposition.

Choose a nonzero $z\in\mathfrak N(R/S)$. This ideal is nonzero
because the generic extension is separable. The first-power
Noether-different result gives
\[
 z\in\mathfrak N(A/S),\qquad
 z\operatorname{Hom}_S(A,S)\subset A\operatorname{Tr}_{A/S}.
                                                        \tag{D.CJ.3}
\]
Here the second inclusion uses the generic trace identification of
the dual module; it also follows directly from the evaluation
description of the Noether different
\citep[Lemma~49.6.7 and complete proof, Tag 0BVT]
{StacksCanonicalJetsB}. Choose a nonzero $c_0\in S\cap zR$;
for example, the constant coefficient of the characteristic
polynomial of multiplication by $z$ on the finite free $S$-module
$R$ has this property by Cayley--Hamilton. Write $c_0=zw$ with
$w\in R$, and put $c=\operatorname{ord}_{\mathfrak q}(c_0)$.

We first prove the following precise scalar estimate:
\[
 a\in H^0(Y,I_b\mathcal O_Y)
                  \quad\Longrightarrow\quad za\in\mathfrak q^bA.
                                                        \tag{D.CJ.4}
\]
Let $v=\operatorname{ord}_{\mathfrak q}$ be the discrete valuation
of $\operatorname{Frac}S$ defined by total order; multiplicativity
holds because $\operatorname{gr}_{\mathfrak q}S=k[X,Y]$ is a domain.
Every extension of $v$ to $\operatorname{Frac}A$ is nonnegative on
$A$, and its valuation ring has a centre on $Y$ by properness
\citep[Lemma~29.43.1, valuative criterion for properness]
{StacksCanonicalJetsB}. At that centre, $a$ is a local linear
combination of $x^b,y^b$. Normalize each extended valuation to
restrict to $v$. Its value on $a$ is therefore at least $b$.

Choose an $S$-basis $e_1,\ldots,e_d$ of $A$, with generic
trace-dual basis $e_1^*,\ldots,e_d^*$. By (D.CJ.3),
$ze_i^*\in A$. Every conjugate of $aze_i^*$ has valuation at
least $b$, so
\[
         \operatorname{Tr}_{A/S}(aze_i^*)\in\mathfrak q^bS.
\]
Indeed its trace belongs to $S$ by integrality, and has total
order at least $b$. Expanding $za$ in the basis $e_i$ proves
(D.CJ.4). The same argument proves
$z\overline{\mathfrak q^bA}\subset\mathfrak q^bA$, but that
integral-closure reformulation is not needed below.

Multiplication by $z$ on $M$ factors through a finite free
$A$-module. Pulling back and taking reflexive hulls gives actual
maps of vector bundles
\[
                 F\xrightarrow{\alpha}\mathcal O_Y^a
                     \xrightarrow{\beta}F,
                 \qquad\beta\alpha=z\operatorname{id}_F.
                                                        \tag{D.CJ.5}
\]
These extensions can be checked at codimension-one discrete
valuation rings and then extended by reflexivity on regular $Y$.
For $m\in K_b$, all coordinates of $\alpha(m)$ belong to
$H^0(Y,I_b\mathcal O_Y)$. Equation (D.CJ.4) consequently gives
\[
 z^2m=\beta(z\alpha(m))\in\mathfrak q^bM,
          \qquad c_0^2m\in\mathfrak q^bM.
\]
The module $M$ is maximal Cohen--Macaulay, hence free over the
regular local ring $S$. Total order in $S$ is additive on
products, so
\[
       (\mathfrak q^bM:c_0^2)=\mathfrak q^{b-2c}M
                   \quad(b\ge2c).
\]
This proves (D.CJ.2). Finally
$\mathfrak q^{2\ell-1}\subset(x^\ell,y^\ell)$ proves the
last assertion. Notice that the valuation argument used the
actual ideal $I_b\mathcal O_Y$; no principalization of that
ideal on $Y$ was assumed.
\end{proof}

\begin{proposition}[Proved: exact two-level recovery of canonical module jets]
\label{D:canonical-module-two-level-recovery}
In the setting of Lemma~\ref{D:canonical-pullback-uniform-jet-loss},
put $r=\operatorname{rank}_A M$ and define the proper parameter-ideal
thickenings and their spaces of sections by
\[
 T_b=V(I_b\mathcal O_Y),\qquad
 V_b(M)=H^0(T_b,F|_{T_b}),\qquad
 W_b(M)=\operatorname{im}\bigl(V_{b+1}(M)\longrightarrow V_b(M)\bigr).
                                                        \tag{D.CJ.6}
\]
The thickenings $T_b$ need not be Cartier or Cohen--Macaulay.
There are natural $A$-module identifications
\[
             W_b(M)=M/K_b,
 \qquad
             W_b(M)/I_\ell W_b(M)\simeq M/I_\ell M
                  \quad(b\ge2\ell-1+2c).
                                                        \tag{D.CJ.7}
\]
They are functorial for homomorphisms of the original reflexive
modules. Thus canonical module jets are recovered from two fixed
parameter-ideal thickenings, with a level independent of the cover
degree and module rank.

Write $d_0=\operatorname{rank}_S R$ and choose $c_R$ such that
$p_g(A)\le c_Rn$, as in
Proposition~\ref{D:cover-uniform-canonical-defect}. Then
\[
 \begin{split}
 h^0(T_b,F|_{T_b})&\le (d_0b^2+2c_R)nr,\\
 h^1(T_b,F|_{T_b})&\le c_Rnr.
 \end{split}                                             \tag{D.CJ.8}
\]
These assertions apply to all actual reflexive modules, without
a Frobenius-periodicity or exceptional-energy hypothesis.
They do not assert effectiveness of an arbitrary family of bundles
on $T_{b+1}$, nor preservation of these images under arbitrary
base change in a parameter scheme.
\end{proposition}

\begin{proof}
First $H^0(Y,F)=M$. One inclusion comes from pullback. Conversely
a section restricts to a section of the bundle associated with
$M$ on the puncture, and Hartogs for the reflexive module $M$
identifies that space of sections with $M$.

We record a useful cohomology estimate with its support argument.
The natural generating sections of $M$ generate $F$ at each
exceptional generic point. Since $k$ is infinite, choose $r$
constant linear combinations whose determinant is a unit at all
those finitely many discrete valuation rings. They give an
injection
\[
                   0\longrightarrow\mathcal O_Y^r
                       \longrightarrow F\longrightarrow Q_F
                       \longrightarrow0.
                                                        \tag{D.CJ.9}
\]
The support of $Q_F$ contains no exceptional component. Its
proper map to $\operatorname{Spec}A$ therefore has finite fibres:
the closed fibre is zero-dimensional, and the resolution is an
isomorphism elsewhere. The support is finite over $A$
\citep[Lemma~37.44.1 and complete proof, Tag 02LS]
{StacksCanonicalJetsB}. Consequently $H^1(Y,Q_F)=0$ and
\[
 H^1(Y,\mathcal O_Y)^r\twoheadrightarrow H^1(Y,F),
 \qquad h^1(Y,F)\le r p_g(A),\qquad JH^1(Y,F)=0.
                                                        \tag{D.CJ.10}
\]
If the exceptional fibre is empty, the same assertion follows
directly from affineness; this also covers the regular case.

Let $D$ be the greatest common divisor of the principal divisors
of $x$ and $y$ on regular $Y$. It is exceptional, since $x,y$
form a parameter ideal. Local unique factorization gives exact
sequences
\[
 0\longrightarrow F(bD)\xrightarrow{\,(-y^b,x^b)\,}
     F^{\oplus2}\xrightarrow{\,(x^b,y^b)\,}
     I_b\mathcal O_Y\cdot F\longrightarrow0.
                                                        \tag{D.CJ.11}
\]
The first arrow is interpreted on the rational sections of
$F(bD)$; its displayed products are regular. Every coherent
sheaf on $Y$ has $H^2=0$, since $\pi$ is proper over an affine
base with fibres of dimension at most one. Hence
\[
          H^1(Y,F)^{\oplus2}\twoheadrightarrow
                      H^1(Y,I_b\mathcal O_Y\cdot F).
                                                        \tag{D.CJ.12}
\]
In the comparison of (D.CJ.11) for $b+1$ and $b$, the map on
the middle terms is $\operatorname{diag}(x,y)$, the map on
the right is inclusion, and the map on the left is
multiplication by $xy$. These maps commute with the displayed
arrows. Since $x,y\in J$, (D.CJ.10)--(D.CJ.12) imply
\[
 H^1(Y,I_{b+1}\mathcal O_Y\cdot F)
                  \longrightarrow H^1(Y,I_b\mathcal O_Y\cdot F)
                        \quad\text{is zero}.
                                                        \tag{D.CJ.13}
\]

Apply cohomology to
$0\to I_b\mathcal O_Y\cdot F\to F\to F|_{T_b}\to0$.
The image of $M$ in $V_b(M)$ is the kernel of its connecting
map to the right-hand group of (D.CJ.13). Naturality and
(D.CJ.13) show that every section from $V_{b+1}(M)$ has zero
connecting image at level $b$. Conversely every section from
$M$ restricts through level $b+1$. It follows that
\[
             W_b(M)=\operatorname{im}(M\to V_b(M))=M/K_b.
\]
The saturation estimate $K_b\subset I_\ell M$ proves the
second identification in (D.CJ.7). All maps in this argument
come from restriction and multiplication by fixed functions,
so the identifications are functorial for actual module maps.

Finally $K_b\supset I_bM$, and $M$ is $S$-free of rank
$d_0nr$. Thus
\[
 \dim_k(M/K_b)\le\dim_k(M/I_bM)=d_0nrb^2.
\]
Equations (D.CJ.10)--(D.CJ.12) bound the possible cokernel of
$M\to V_b(M)$ by $2r p_g(A)$. The same restriction sequence
and $H^2(Y,I_b\mathcal O_Y\cdot F)=0$ give a surjection
$H^1(Y,F)\to H^1(T_b,F|_{T_b})$. These are exactly (D.CJ.8).
\end{proof}

\begin{remark}[The remaining all-family realization problem]
The proposition is an exact pointwise recovery, rather than
an inference from annihilation to stabilization. In particular
it specifies the transition whose image is the correct one.
The schemes $T_b$ are parameter-ideal thickenings and can have
embedded zero-dimensional structure; replacing them silently
by reduced curves or Cartier cycles would change the statement.
To obtain a quantitative hereditary representation atlas one
still has to construct, over arbitrary Artinian parameter
algebras, effective object and arrow diagrams whose section
images realize (D.CJ.7), and control their equation degrees.
Neither a base-change assertion for these images nor a bound
on the resulting number of representation classes is proved here.
\end{remark}
